% TOP_PROBLEM: {"problemId": "problem.calderon-uniqueness-for-bounded-measurable-isotropic-conductivities-in-higher-dimensions", "canonicalTitle": "Calderón uniqueness for bounded measurable real isotropic conductivities", "releaseRank": 314, "primaryDomain": "analysis_pde", "primaryDomainLabel": "Analysis & PDE", "displayStatus": "open_with_solved_subcases", "releaseStatus": "open_with_solved_subcases"}
% !TeX program = lualatex
% Definition and sources only; this file is not an LLM solution attempt.
\documentclass[11pt]{article}
\usepackage[margin=25mm]{geometry}
\usepackage{fontspec}
\setmainfont{Latin Modern Roman}
\usepackage{amsmath,amssymb,mathtools}
\usepackage{unicode-math}
\setmathfont{Latin Modern Math}
\usepackage{xurl,hyperref}
\hypersetup{colorlinks=true,urlcolor=blue}
\setcounter{secnumdepth}{0}
\setlength{\parindent}{0pt}
\setlength{\parskip}{5pt}
\setlength{\emergencystretch}{3em}
\begin{document}
\section*{\texorpdfstring{Calderón uniqueness for bounded measurable real isotropic conductivities}{Calderón uniqueness for bounded measurable real isotropic conductivities}}\label{314}
\subsection{Definitions and mathematical statement}
Let $n\ge3$ and let $\Omega\subset{\mathbb{R}}^n$ be a bounded connected Lipschitz domain. For a real scalar conductivity $\gamma\in L^\infty(\Omega)$ with $0<c\le\gamma\le C<\infty$, the weak solution with boundary value $f\in H^{1/2}(\partial\Omega)$ satisfies $\nabla\cdot(\gamma\nabla u)=0$. Its Dirichlet-to-Neumann map is the bounded operator determined by
\[
 \langle\Lambda_\gamma f,g\rangle
 =\int_\Omega\gamma\nabla u\cdot\nabla v{\,\mathrm{d}} x,
 \qquad v|_{\partial\Omega}=g.
\]
The value is independent of the chosen $H^1$ extension $v$. The uniqueness question is whether $\Lambda_{\gamma_1}=\Lambda_{\gamma_2}$ implies $\gamma_1=\gamma_2$ almost everywhere. The hypotheses impose no derivative regularity, smallness, or prior agreement near the boundary. This scalar question is distinct from anisotropic uniqueness modulo boundary-fixing changes of coordinates.
\par\smallskip\textit{Formulation and concept references:} \href{https://arxiv.org/pdf/1803.06931v2}{[S1]}.

\subsection{Short English statement}
Do all boundary voltage-to-current measurements uniquely determine a bounded measurable positive scalar conductivity in dimension at least three?
\subsection{Sources}
\begin{itemize}
\item[S1]Matteo Santacesaria, Note on Calderón's inverse problem for measurable conductivities, Inverse Problems and Imaging13(1)(2019),149–157; arXiv1803.06931v2.
\url{https://arxiv.org/pdf/1803.06931v2}.
\textit{Formulation source.}
\end{itemize}
\end{document}
